\chapter{KẾT LUẬN VÀ HƯỚNG PHÁT TRIỂN}
\label{chuong5}
\label{ch:ket_luan}

\vspace{-1.5cm}

\begin{quotation}
    \noindent
    \small\itshape
    Chương này tổng kết kết quả thực nghiệm của khóa luận, chuyển phần kiểm định chéo sang phần thảo luận cuối cùng, trả lời ba câu hỏi nghiên cứu đã nêu ở Chương~\ref{ch:tong_quan}, và nêu rõ giới hạn của đề tài. Các hướng phát triển được giới hạn ở những thí nghiệm và mở rộng dữ liệu trực tiếp xuất phát từ kết quả hiện tại, tránh mô tả các ứng dụng quá rộng ngoài phạm vi khóa luận.
\end{quotation}
\vspace{1em}

\section{Thảo luận kết quả thực nghiệm}
\label{sec:thao_luan_ket_qua}

\subsection{Kết quả theo bộ dữ liệu}

Kết quả ở Chương~\ref{ch:ket_qua_thuc_nghiem} cho thấy AutoShotV2 đạt hiệu quả tốt nhất trên SHOT và BBC trong cấu hình chính. Trên SHOT, mô hình đạt F1 = \ASVTwoShotFOne{}, cao hơn AutoShot gốc khoảng $+\ASVTwoVsAutoShotShotPP\%$ và cao hơn TransNetV2 khoảng $+\ASVTwoVsTransNetShotPP\%$. Điều này phù hợp với mục tiêu của khóa luận: giữ lại mạng xương sống AutoShot và cải thiện tầng phân loại, nhãn giám sát và hậu xử lý xác suất.

Với BBC, AutoShotV2 đạt F1 = \ASVTwoBBCFOne, cho thấy cấu hình cân bằng vẫn hoạt động tốt trên video tài liệu dài. Tuy nhiên, trên ClipShots, AutoShotV2 chỉ đạt F1 = \ASVTwoClipFOne, thấp hơn AutoShot theo báo cáo gốc $3.4\%$. Vì vậy, kết quả không nên được hiểu là mô hình đã tổng quát hóa tốt trên mọi miền video; kết luận hợp lý hơn là AutoShotV2 cải thiện rõ trên SHOT, giữ mức tốt trên BBC, nhưng còn yếu khi gặp miền dữ liệu đa dạng hơn như ClipShots.

\subsection{Phân tích suy giảm trên ClipShots}

Bảng~\ref{tab:clipshots_breakdown} phân tách kết quả AutoShotV2 trên ClipShots theo loại chuyển cảnh. Bảng này dùng để phân tích lỗi theo loại chuyển cảnh, không thay thế kết quả F1 chính ở Bảng~\ref{tab:autoshotv2_prf}.

\begin{table}[H]
\centering
\caption[Breakdown theo loại chuyển cảnh trên ClipShots]{Breakdown Precision/Recall/F1 theo loại chuyển cảnh trên ClipShots test (500 video, ngưỡng = 0.10).}
\label{tab:clipshots_breakdown}
\small
\begin{tabular}{lccc}
\toprule
\textbf{Loại chuyển cảnh} & \textbf{Precision} & \textbf{Recall} & \textbf{F1} \\
\midrule
Cut (chuyển cảnh đột ngột) & 0.5826 & 0.8860 & 0.7030 \\
Gradual (chuyển cảnh dần) & 0.6456 & 0.0442 & 0.0828 \\
\textbf{Tổng}              & \textbf{0.7477} & \textbf{0.7897} & \textbf{0.7681} \\
\bottomrule
\end{tabular}
\end{table}

Điểm yếu lớn nhất nằm ở chuyển cảnh dần: Recall của nhóm \textit{gradual} chỉ đạt $4.42\%$. Điều này cho thấy mô hình gần như bỏ sót phần lớn các chuyển cảnh dần trong ClipShots. Trong phạm vi khóa luận, nguyên nhân chỉ được phân tích ở mức giả thuyết: nhãn many-hot $\pm 1$ frame có thể chưa phù hợp với ranh giới sắc của ClipShots, temperature scaling được tối ưu trên SHOT có thể lệch khi sang miền khác, và cấu hình Focal Loss hiện tại chưa chắc tối ưu cho phân phối nhãn của ClipShots.

\subsection{Giới hạn diễn giải}

Các nhận định trên được rút ra từ F1, Precision, Recall và một số bảng phân rã lỗi hiện có. Khóa luận chưa thực hiện đầy đủ thí nghiệm đối chứng để tách riêng từng yếu tố như many-hot label, temperature scaling hay Focal Loss trên từng bộ dữ liệu. Vì vậy, các phân tích về ClipShots nên được hiểu là cơ sở giải thích và định hướng cải thiện, không phải khẳng định nhân quả tuyệt đối.

\section{Kiểm định chéo bổ sung theo từng bộ dữ liệu}
\label{sec:calibration}

Phần này được chuyển từ Chương~\ref{ch:ket_qua_thuc_nghiem} sang Chương~\ref{ch:ket_luan} vì vai trò chính của nó là kiểm chứng lại xu hướng kết quả, không phải là bảng kết quả chính. Kiểm định chéo 5 phần (5-fold cross-validation) giúp hạn chế việc diễn giải quá mức từ một lần chọn ngưỡng trực tiếp trên tập kiểm tra.

\begin{table}[H]
\centering
\caption[F1 kiểm định chéo theo từng bộ dữ liệu]{F1 kiểm định chéo 5 phần theo từng bộ dữ liệu. Số đậm là giá trị tốt nhất theo cột.}
\label{tab:calibration_cv}
\small
\begin{tabular}{lccc}
\toprule
\textbf{Mô hình} & \textbf{SHOT} & \textbf{BBC} & \textbf{ClipShots} \\
\midrule
\CalibrationCVRows
\bottomrule
\end{tabular}
\end{table}

Kết quả kiểm định chéo củng cố ba điểm chính. Thứ nhất, baseline $A0$ tăng rõ trên ClipShots, từ $0.7649$ lên $0.7881$ ($+2.32\%$), cho thấy ClipShots nhạy với cách chọn ngưỡng. Thứ hai, Pha 2 giúp SHOT ($+1.0\%$) và BBC ($+0.9\%$) nhưng vẫn kém $A0$ trên ClipShots ($-3.1\%$). Thứ ba, mức trần khi chọn trực tiếp trên tập kiểm tra khá gần với số kiểm định chéo, ví dụ ClipShots $A0$ chênh $+0.15\%$, nên xu hướng kết quả là tương đối ổn định.

\section{Kết luận chung}
\label{sec:ket_luan}

\subsection{Trả lời các câu hỏi nghiên cứu}

\textbf{RQ1 -- NAS có giúp cải thiện hiệu năng SBD trên video ngắn so với mạng xương sống cố định?}
Câu trả lời là \textbf{có}. Trên SHOT, AutoShot đạt F1 = \AutoShotOriginalShotFOne, vượt TransNetV2 (F1 = \TransNetShotFOne) khoảng $+4.1\%$. Điều này cho thấy kiến trúc được tìm bằng NAS là điểm khởi đầu mạnh hơn mạng xương sống cố định trong bối cảnh video ngắn.

\textbf{RQ2 -- Các kỹ thuật ở tầng phân loại có cải thiện F1 mà không cần thay đổi mạng xương sống?}
Câu trả lời là \textbf{có}. AutoShotV2 giữ nguyên mạng xương sống AutoShot nhưng bổ sung Focal Loss, nhãn many-hot, Gaussian smoothing và temperature scaling. Trên SHOT, mô hình đạt F1 = \ASVTwoShotFOne{}, cao hơn AutoShot gốc $+\ASVTwoVsAutoShotShotPP\%$.

\textbf{RQ3 -- Mức cải thiện trên SHOT có duy trì khi đánh giá chéo trên BBC và ClipShots?}
Câu trả lời là \textbf{một phần}. AutoShotV2 giữ kết quả tốt trên BBC nhưng suy giảm trên ClipShots. Do đó, cấu hình hiện tại phù hợp với miền gần SHOT hơn là một lời giải tổng quát cho mọi loại video ngắn.

\section{Đóng góp và hạn chế của khóa luận}
\label{sec:han_che}

Trong phạm vi khóa luận này, đóng góp chính không phải là xây dựng một hệ thống sản phẩm hoàn chỉnh, mà là kiểm chứng một hướng cải tiến có kiểm soát trên nền AutoShot:

\begin{enumerate}
    \item Hệ thống hóa bài toán SBD, các nhóm mô hình liên quan và vai trò của AutoShot trong hướng NAS cho video ngắn.
    \item Triển khai và đánh giá AutoShotV2 với các cải tiến ở tầng phân loại và hậu xử lý, trong khi giữ nguyên mạng xương sống AutoShot.
    \item Báo cáo kết quả trên SHOT, BBC và ClipShots, kèm phân tích rõ giới hạn tổng quát hóa trên ClipShots.
\end{enumerate}

Khóa luận cũng có các hạn chế cần nêu rõ. Thứ nhất, quy trình NAS gốc không được chạy lại từ đầu do chi phí tính toán lớn; khóa luận kế thừa kiến trúc và trọng số công khai của AutoShot để tập trung vào tầng phân loại. Thứ hai, các thí nghiệm phân rã chưa đủ rộng để kết luận chắc chắn từng yếu tố gây suy giảm trên ClipShots. Thứ ba, khóa luận mới đánh giá trên SHOT, BBC và ClipShots, nên chưa thể khẳng định mức ổn định khi chuyển sang các tập video ngắn khác hoặc dữ liệu do người dùng tự thu thập.

\section{Hướng phát triển trong phạm vi khóa luận}
\label{sec:huong_phat_trien}

Các hướng phát triển sau bám sát các kết quả đã quan sát được, chủ yếu nhằm kiểm chứng nguyên nhân và cải thiện độ ổn định của AutoShotV2:

\begin{enumerate}
    \item \textbf{Thực hiện thí nghiệm phân rã trên ClipShots.} Cần chạy lại từng biến thể loại bỏ many-hot label, thay đổi temperature scaling, và quét lại các tham số $\gamma, \alpha$ của Focal Loss để xác định yếu tố nào ảnh hưởng mạnh nhất đến suy giảm F1 trên ClipShots.

    \item \textbf{Hiệu chỉnh theo từng miền dữ liệu.} Thay vì dùng một ngưỡng và một hệ số nhiệt độ cố định cho mọi bộ dữ liệu, có thể xây dựng quy trình chọn $\theta$ và $T^*$ trên một tập xác thực nhỏ của miền đích. Hướng này trực tiếp giải quyết khác biệt giữa SHOT, BBC và ClipShots đã thể hiện trong kết quả thực nghiệm.

    \item \textbf{Mở rộng đánh giá trên dữ liệu video ngắn khác.} Một tập đánh giá bổ sung, được gán nhãn theo cùng quy ước với SHOT, sẽ giúp kiểm tra xem kết luận hiện tại có còn giữ khi thay đổi nguồn video, phong cách dựng và tỉ lệ chuyển cảnh dần hay không.
\end{enumerate}

Những hướng trên giữ trọng tâm ở mức nghiên cứu thực nghiệm: kiểm chứng giả thuyết, cải thiện quy trình hiệu chỉnh và mở rộng đánh giá có kiểm soát, thay vì trình bày AutoShotV2 như một hệ thống ứng dụng hoàn chỉnh.

\section{Lời kết}

Khóa luận đã khảo sát và thực nghiệm hướng cải tiến AutoShot cho bài toán phát hiện ranh giới cảnh quay. Kết quả cho thấy AutoShotV2 cải thiện rõ trên SHOT và duy trì kết quả tốt trên BBC, nhưng còn hạn chế trên ClipShots. Đóng góp chính của khóa luận là chỉ ra rằng các cải tiến ở tầng phân loại và hậu xử lý có thể nâng hiệu quả mô hình mà không cần thay đổi mạng xương sống, đồng thời làm rõ giới hạn tổng quát hóa cần được kiểm chứng thêm trong các thí nghiệm tiếp theo.
